\documentclass[11pt,letterpaper]{article}
\usepackage[T1]{fontenc}
\usepackage[utf8]{inputenc}
\usepackage{lmodern,microtype}
\usepackage[margin=0.88in,headheight=15pt]{geometry}
\usepackage{amsmath,amssymb,amsthm}
\usepackage{booktabs,longtable,array,tabularx}
\usepackage{xcolor,enumitem,listings,fancyhdr}
\usepackage{hyperref,xurl}
\definecolor{ink}{HTML}{173447}
\definecolor{accent}{HTML}{226877}
\definecolor{muted}{HTML}{53616A}
\hypersetup{colorlinks=true,linkcolor=ink,urlcolor=accent,citecolor=accent,
  pdftitle={Asymptotic: Incremental Repository Audit},
  pdfauthor={Independent technical review},bookmarksnumbered=true}
\pagestyle{fancy}
\fancyhf{}
\lhead{\small\textcolor{muted}{ASYMPTOTIC \; / \; INCREMENTAL AUDIT}}
\rhead{\small\textcolor{muted}{Snapshot 651f2029}}
\cfoot{\small\thepage}
\setlength{\parindent}{0pt}
\setlength{\parskip}{6pt plus 1pt}
\setlength{\emergencystretch}{3em}
\setlist{nosep,leftmargin=1.4em}
\renewcommand{\arraystretch}{1.17}
\newcolumntype{Y}{>{\raggedright\arraybackslash}X}
\newcommand{\id}[1]{\textbf{\textcolor{ink}{#1}}}
\newcommand{\prop}[1]{\texttt{#1}}
\newcommand{\pin}{651f2029d0b2cd4da9e4dfdf1f4275a124d23b99}
\newcommand{\sourcepath}[1]{\path{#1}}
\newtheorem{proposition}{Proposition}
\lstdefinelanguage{Wolfram}{morekeywords={Module,Block,If,Return,Table,Do,With,
  Get,Check,SeriesTruncate,AsymptoticExpansion,VerificationTest,TestReport,Normal,
  DownValues,RuleDelayed,Missing,True,False,Infinity,Log,Zeta,Abs,Options},
  sensitive=true,morecomment=[s]{(*}{*)},morestring=[b]"}
\lstset{basicstyle=\ttfamily\footnotesize,columns=fullflexible,keepspaces=true,
  breaklines=true,breakatwhitespace=false,showstringspaces=false,
  keywordstyle=\color{ink}\bfseries,commentstyle=\color{muted},
  stringstyle=\color{accent},frame=single,framerule=0.3pt,rulecolor=\color{muted!35},
  xleftmargin=0.4em,xrightmargin=0.4em,aboveskip=9pt,belowskip=9pt}

\begin{document}
\begin{titlepage}
\thispagestyle{empty}
\vspace*{0.45in}
{\Large\textcolor{accent}{SOFTWARE, MATHEMATICS, AND VALIDATION}}\par
\vspace{0.22in}
{\Huge\bfseries\textcolor{ink}{Asymptotic}}\par
\vspace{0.08in}
{\LARGE Incremental repository audit}\par
\vspace{0.2in}
{\large Remainder-certificate provenance, affine-recognition complexity,\par
and consistency of acceptance evidence}\par
\vspace{0.32in}
\rule{\linewidth}{0.6pt}
\textbf{Reviewed repository}\quad VladimirReshetnikov/Asymptotic\par
\textbf{Pinned revision}\quad\texttt{651f2029d0b2cd4da9e4dfdf1f4275a124d23b99}\par
\textbf{Review date}\quad September 10, 2026\par
\rule{\linewidth}{0.6pt}
\vspace{0.2in}
\textbf{Abstract}\par
This report examines the current implementation against the repository's existing
four-wave review corpus, rather than presenting its established backlog as new
work. Three additional finding groups survive that exclusion process. First, a
second, no-op truncation retains a transported quantitative remainder contract but
loses the companion expression that the contract explicitly names. A focused
public Wolfram-kernel witness confirms this inconsistency; it does not show an
incorrect error bound. Second, the newly added exact-affine fast path repeatedly
reclassifies overlapping nonlinear subtrees. For a Horner expression with
$4n+1$ tree nodes, the inspected control flow performs exactly $4n(n+1)$ affine
recognizer visits. A separate exact-rational model reproduces the count and
preserves interval endpoints under request-local classification reuse. Third,
the new Mathics acceptance checker accepts synthetic receipts with contradictory
entry-point or source-stability metadata, despite checking their source maps and
protocol records successfully.

The accompanying artifacts include a source-matched candidate truncation patch,
a candidate receipt guard, an independent interval prototype, runnable Python
checks, and Wolfram regression specifications. The executed Python run has 34
passing test cases. Native observations, synthetic records, source proofs, and
unrun acceptance specifications are kept separate throughout. This is a focused
incremental audit, not a claim of exhaustive defect discovery or full-suite
acceptance.
\vfill
{\small\textcolor{muted}{The snapshot is the unit of analysis. Later changes to
\texttt{main} are outside the claims in this report. No remote repository changes
were made.}}
\end{titlepage}

\begingroup
\footnotesize
\setlength{\parskip}{0pt}
\tableofcontents
\endgroup
\newpage

\section{Executive assessment}
\label{sec:executive}

The most useful new work in this snapshot is not another general request for
better caching, richer transseries, native compatibility, or stronger certificates.
Those directions already appear in the existing reviews. The additional defects
identified here occur at recently changed boundaries: a certificate copied after
truncation, a speculative recognizer inserted before interval evaluation, and a
postprocessor that consolidates test receipts. Each boundary can appear locally
reasonable while violating a compositional property one level above it.

\begin{table}[htbp]
\small
\begin{tabularx}{\linewidth}{@{}p{0.06\linewidth}p{0.17\linewidth}Yp{0.23\linewidth}@{}}
\toprule
ID & Classification & New observation & Evidence obtained \\
\midrule
N01 & Medium, metadata integrity & A no-op truncation retains a
\prop{TransportedThroughTruncation} contract but drops
\prop{TruncationDiscardedPart}. & Pinned source and one successful public native witness. \\
N02 & Medium, performance & Repeated failed affine recognition costs
$\Theta(n^2)$ structural visits on an $O(n)$ Horner tree. & Source recurrence,
exact model, endpoint-preserving prototype checks. \\
N03 & Medium, validation integrity & The new receipt checker accepts selected-entry
and recorded-drift contradictions. & Executed upstream-function excerpt against
explicitly synthetic records; candidate-guard controls. \\
\bottomrule
\end{tabularx}
\caption{Finding groups retained after comparison with the existing review baseline.
Severity describes the demonstrated consequence, not a hypothetical worst case.}
\end{table}

None of these observations establishes a newly incorrect asymptotic coefficient,
a false rational interval enclosure, or an actually corrupted historical CI
receipt. That distinction matters. N01's stored upper bound remains valid; N02
is an algorithmic cost defect in a helper path; N03 is a demonstrated acceptance
blind spot on artificial records. Inflating these conclusions would obscure both
their real significance and the work still required to close them.

The immediate actions are correspondingly small and testable. Preserve the
existing discard companion on the no-op branch. Compute affine classifications
once per interval-evaluation request without weakening the exact-affine
translation fix. Make the receipt checker reject a selected source that is not
the intended fingerprinted entry, and reject a non-null record of observed
source drift. Broader changes should follow only after these exact invariants
have focused acceptance.

\section{Snapshot, scope, and method}
\label{sec:scope}

\subsection{What was reviewed}

The reviewed revision is \texttt{\pin}. Its commit records the article's updated
account of affine ranges and truncation-bound transport. The maintained package
is named \prop{AsymptoticAnalysis}; the public inverse constructor remains
\prop{AsymptoticInverse}. The repository contains modular kernel sources and a
generated standalone entry. The latter was successfully imported as a complete
685,212-character text and loaded through a string stream during this review.
This observation concerns the pinned standalone file, not every possible loading
mode. Source and development references are commit-pinned in the bibliography.
\cite{repo,core,status}

The source examination covered the ordinary power-log representation and
precision algebra, forward and inverse dispatch, explicit series operations,
Lambert and reciprocal-log handling, Gamma/Barnes forward adapters, convergent
Dirichlet and Lerch expansions, rational interval certificates, several Mathics
compatibility modules, native-result boundaries described by the implementation
register, and the new acceptance checker. Inspection was particularly detailed
around the paths relevant to the retained findings. The audit does not claim that
every line of every kernel module, test, vendor document, and external report was
independently reread.

The exclusion baseline contains 36 review packages in four waves. The maintained
register reports 231 attributed entries before overlap consolidation. That number
is neither a count of distinct defects nor a count of outstanding problems.
The register and the two later-wave intakes were used as the principal exclusion
map, with targeted checks of the original nearby reports. In particular, report
14's original affine-rounding and quantitative-bound findings were compared with
their current repairs. Report 36's ledger was checked to distinguish its runner
and aggregation concerns from the new downstream checker examined here.
\cite{reviews,status,w3,w4,r14,r36}

\subsection{Evidence categories}

Four categories are deliberately not combined into a single passing total.
\emph{Source-proved behavior} is a consequence of the inspected code, with the
relevant admission and control-flow assumptions stated. \emph{Native observation}
means that a successful official-kernel tool call returned the reported public
result. \emph{Independent-model evidence} executes an explicitly separate model
of an algorithm or mathematical identity. \emph{Synthetic-validator evidence}
executes validation functions against constructed records; it says nothing about
whether Mathics actually ran those cases.

The bundled Wolfram files contain further acceptance specifications and helper
instrumentation. Those files were not successfully executed as a suite in this
review. Subsequent attempts to use the native service returned network/502
failures, including a simple version query. Such attempts contribute no test
results. Earlier partial direct-URL loading attempts likewise do not count as
successful package loads. The successful string-stream load and public N01
witness are recorded separately in \path{evidence/native-witness.json}.

\subsection{How novelty was decided}

A finding was excluded when its mechanism and proposed obligation were already
represented in the prior corpus, even if it remained open. A changed witness
alone was not treated as novelty. Conversely, a defect introduced by a repair can
be new when it concerns a newly added field, call path, or verifier obligation
that did not exist in the reviewed baseline. The retained findings meet this
narrower criterion; their nearest prior items remain explicitly attributed.

This method avoids two common distortions. A fixed historical issue is not
reclassified as a current defect merely because its original report is still in
the repository. A general old recommendation is not renamed and presented as a
new architectural proposal. The development plan in Section~\ref{sec:plan}
therefore concerns only the specific new obligations demonstrated below.

\section{Architecture and the assessed boundaries}
\label{sec:architecture}

The ordinary representation separates a positive local coordinate, a finite
jet, and a remainder class. Explicit calculus uses an expression of the form
\[
  b(x)+p(x)\left(\sum_j w(x)^{\beta_j} C_j(\log w(x))+R(w(x))\right),
  \qquad w(x)\to0^+,
\]
where $b$ is an offset and $p$ is an exact carrier. The jet records retained rows,
a power frontier, and a logarithmic degree. This is substantially more structured
than a bare finite expression. It explains why preservation of metadata is part
of correctness of the object API, even when a displayed approximation is
unchanged.\cite{core,operations}

The inspected arithmetic already accounts for precision loss in multiplication,
complete coefficient rows, and logarithmic degrees. The current register also
records focused repairs for real-coefficient admission, retained assumptions,
semantic equality of exponents, native-tail import/export, and separate Fourier
termination. These repairs are the baseline of this report, not additional
findings. Their recorded acceptance counts remain scoped to their own snapshots
and selected suites.\cite{status,w3}

\begin{longtable}{@{}p{0.25\linewidth}p{0.38\linewidth}p{0.30\linewidth}@{}}
\toprule
Boundary & Role in the implementation & Outcome of this audit \\
\midrule
\endfirsthead
\toprule
Boundary & Role in the implementation & Outcome of this audit \\
\midrule
\endhead
Ordinary jet algebra & Retained power-log blocks and inherited precision are
combined before public result construction. & No additional coefficient
counterexample established in the inspected paths. \\
Explicit calculus & Transports offsets, carriers, jets, assumptions, domains,
and recipes. & N01 concerns the new quantitative-bound companion-copy branch. \\
Lambert and reciprocal-log paths & Keep exact cores or carriers and use a
separate small inverse-logarithmic coordinate. & Existing closure, domain,
and cutoff concerns were excluded rather than rediscovered. \\
Gamma/Barnes forward normalization & Positive special-function factors are
combined in the real logarithmic domain, with bounded simplification. & No new
identity error established in the inspected normalization paths. \\
Dirichlet/Lerch providers & Supply explicit coefficient and quantitative-tail
metadata from convergent defining sums and moment bounds. & A Zeta result gives
a transparent N01 witness; the provider's bound is not the defect. \\
Exact certificate engine & Uses rational interval endpoints and outward
rounding; the new affine path delays rounding until an exact range is found. &
N02 identifies repeated recognition work, not unsound rounding. \\
Mathics adapters & Supply bounded package-local compatibility operations and
avoid modifying the global System implementations. & Previously reported
semantic and resource limitations remain excluded. \\
Acceptance postprocessor & Matches recorded inputs, protocol output, cases,
versions, and coverage against a reference. & N03 finds missing consistency
checks between fields that the postprocessor otherwise validates separately. \\
\bottomrule
\end{longtable}

The separate mathematical scales are intentional, not themselves evidence of a
missing universal representation. Likewise, a native formal expansion is not
silently treated here as an analytic remainder theorem. Those distinctions are
already explicit in the maintained register and related notes. This review does
not count the known native-superset, Mathics-completeness, or vendored-corpus goals
as newly discovered unimplemented features.\cite{status,w3,w4}

\section{Novelty crosswalk}
\label{sec:novelty}

\begin{table}[htbp]
\small
\begin{tabularx}{\linewidth}{@{}p{0.07\linewidth}p{0.22\linewidth}Y@{}}
\toprule
New ID & Nearest prior item & What is different here \\
\midrule
N01 & R14 N05; maintained C22 & The original quantitative bound is now preserved.
The newly introduced \prop{TruncationDiscardedPart} field is subsequently omitted
by a no-op metadata copy, leaving its retained contract incomplete. \\
N02 & R14 N02; maintained C20; general P08 & Exact affine evaluation before rounding
is now implemented. Its recursive placement causes a new, explicitly derived
quadratic recognizer count on nested nonlinear expressions. \\
N03 & W4-09, W4-12, W4-14; R36 N06 & The downstream checker is a later implementation.
The counterexamples concern contradictions in its saved-record validation, not
rediscovery of source-freezing, package-load, or process-aggregation gaps. \\
\bottomrule
\end{tabularx}
\caption{Attribution without repeating the old recommendation as a new finding.}
\end{table}

The relevant distinction for N01 is especially narrow. Report 14 already showed
that quantitative tail metadata could be lost by truncation and proposed
transport through discarded known terms. The current code implements that
transport. The present report neither repeats the old missing-bound witness nor
asks again for that general feature. It demonstrates that the repair's new
companion field is not included in the repair's own no-op copy list.
\cite{r14,operations,status}

N02 also preserves the earlier repair's purpose. Report 14's affine example shows
why rounding a huge translated constant before cancellation can destroy a useful
certificate. The new performance proposal must retain exact affine cancellation;
removing the fast path would regress the old finding. The new analysis concerns
how often the exact classifier is invoked after that fast path was inserted.
\cite{r14,certificates}

For N03, repository history identifies the addition of
\path{validation/check_mathics_acceptance.py} at commit
\texttt{e10ef1b71f5b55080c62f52c3ce9d3d6e039fc99}, whose message describes
verifying and recording complete focused Mathics acceptance. The inspected
validator therefore has a concrete implementation and a separate proof obligation.
The synthetic counterexamples do not allege a defect in the specific acceptance
record added by that commit.\cite{checker,checkercommit}

\section{N01: a retained bound contract loses its discard companion}
\label{sec:n01}

\subsection{The public witness}

Use the package backend explicitly, so native fallback is not an alternative
explanation of the result:

\begin{lstlisting}[language=Wolfram]
s = AsymptoticAnalysis`AsymptoticExpansion[
  Zeta[x], {x, Infinity, Log[4]}, "Backend" -> "Package"];
t = AsymptoticAnalysis`SeriesTruncate[s, Log[2]];
u = AsymptoticAnalysis`SeriesTruncate[t, Log[2]];

{t["TruncationDiscardedPart"],
 u["TruncationDiscardedPart"],
 t["AbsoluteRemainderBound"] === u["AbsoluteRemainderBound"],
 u["ForwardRemainderContract"]["Type"], Normal[t], Normal[u]}
\end{lstlisting}

The successful native tool call returned three \prop{GeneralizedSeries} objects.
After alpha-renaming its generated variable to $x$, the relevant observation is
\begin{lstlisting}[language=Wolfram]
{2^(-x) + 3^(-x),
 Missing["KeyAbsent", "TruncationDiscardedPart"],
 True,
 "TransportedThroughTruncation",
 1, 1}
\end{lstlisting}

This call also emitted \prop{Power::infy} and \prop{Greater::nord} diagnostics.
Their origin was not isolated, and they are not counted as a separate finding.
The stored-field comparison, object heads, and unchanged expressions are the
specific observations used here. The service had separately identified itself
as Wolfram 15.0.0 on Linux; that version was not re-emitted inside this witness's
output. The evidence file preserves these qualifications.

\subsection{Why the upper bound remains correct}

For real $x>1$, the provider used by this example stores the first three
Dirichlet terms and a positive omitted tail:
\[
 \zeta(x)=1+2^{-x}+3^{-x}+R_4(x),
 \qquad
 0<R_4(x)\le B_4(x):=4^{-x}\left(1+\frac{4}{x-1}\right).
\]
The bound follows from the first omitted term plus the integral of $t^{-x}$ from
$4$ to infinity. It is also the formula implemented by the inspected provider.
\cite{dirichlet}

The first truncation discards
\[
 \Delta(x)=2^{-x}+3^{-x}
\]
and retains the approximation $1$. The current transport code correctly forms
\[
 B_t(x)=B_4(x)+|\Delta(x)|.
\]
Consequently $|\zeta(x)-1|\le B_t(x)$ under the retained conditions. The second
truncation discards no additional row and keeps the same $B_t$. Nothing in this
witness makes $B_t$ false or changes the power-log remainder. The defect is that
the second result still advertises a contract whose statement refers to
\prop{TruncationDiscardedPart}, although that property is now missing.

\subsection{The exact source mechanism}

In \path{src/Kernel/SeriesOperations.wl},
\prop{seriesTransportRemainderBound} implements the transport. The nontrivial branch computes the removed
rows, constructs their expression, and attaches both a bound and a companion.
The following excerpt omits unrelated fields and abbreviates the statement:
\begin{lstlisting}[language=Wolfram]
"AbsoluteRemainderBound" ->
  data["AbsoluteRemainderBound"] + Abs[discarded],
"TruncationDiscardedPart" -> discarded,
"ForwardRemainderContract" -> <|
  "Type" -> "TransportedThroughTruncation",
  "OriginalAbsoluteRemainderBound" -> data["AbsoluteRemainderBound"],
  "OriginalContract" -> Lookup[data, "ForwardRemainderContract", ...],
  "Statement" -> "... plus Abs[TruncationDiscardedPart] ..."
|>
\end{lstlisting}

When no rows are removed, however, the helper builds a fresh result and copies
only a fixed list of optional keys. That list contains
\prop{ForwardRemainderContract} and the quantitative bound, but not
\prop{TruncationDiscardedPart}. The previous companion is neither regenerated nor
copied. The public truncation wrapper invokes this helper after constructing the
new result, so the omission is directly reachable.\cite{operations}

The prerequisite is an already transported object. Applying a no-op to an
original Zeta expansion need not have a discard companion at all. A regression
must therefore exercise the two-stage sequence, rather than merely checking
whether a single truncation preserves \prop{AbsoluteRemainderBound}.

\subsection{A minimal repair and its invariant}

Add \prop{TruncationDiscardedPart} to the existing no-op copy list. The guard that
copies only keys actually present should remain unchanged. This preserves the
companion when it exists and does not invent a history field for an original
untransported result. The nontrivial branch should continue computing its newly
discarded part rather than reusing a previous one.

The desired invariant is a scoped form of idempotence. Let $T_h$ be truncation at
cutoff $h$, and suppose a second application removes no more rows. Then the finite
expression and remainder should remain unchanged, and any copied transported-bound
contract should retain the companion on which its statement depends. This does
not require byte-for-byte equality of whole objects: a new recipe wrapper or
other provenance may legitimately differ.

The supplied patcher performs this one-list change only when the pinned-source
text matches exactly once. It defaults to showing a diff and writes only with
\texttt{--write}. It changes the modular source; the generated standalone must be
rebuilt by the repository's own builder. Python fixtures check matching,
line-ending preservation, and refusal of unmatched or repeated patches. These
fixtures validate the patcher's operation, not the patched Wolfram package.

The six supplied MUnit specifications check companion preservation, unchanged
bounds, unchanged expression/remainder, repeated no-ops, absence of fabricated
history on an original object, and retained contract type. They have not been
executed as a native suite in this review. Focused native acceptance of the patch
is therefore still required.

\section{N02: repeated affine recognition is quadratic on Horner trees}
\label{sec:n02}

\subsection{The new fast path is mathematically useful}

The certificate engine represents an enclosure by two exact rational endpoints
and rounds outward to a dyadic grid. Its new affine recognizer returns a pair
$(a,b)$ for an expression equal to $ax+b$, provided that the expression is built
from rational constants, the source symbol, addition, and multiplication of
admitted affine factors. A product is refused when two accumulated factors are
both nonconstant. An admitted affine expression is evaluated on an interval $I$
by the exact range
\[
 \operatorname{range}(ax+b,I)=
 \bigl[\min_{z\in\partial I}(az+b),\,
       \max_{z\in\partial I}(az+b)\bigr],
\]
followed by one outward rounding step.\cite{certificates}

This fixes a genuine conditioning problem in expression evaluation: at a huge
rational translation, rounding $x$ and the translation independently can lose
cancellation that is exact in the original affine function. The performance
problem discussed here is not an objection to that mathematics.

\subsection{Where repeated work enters}

\prop{certEnclose} first tries \prop{certAffineRange} on every nontrivial
\prop{Plus} or \prop{Times} node. If the recognizer fails, ordinary interval
evaluation recurses into the children. Those children then try the same affine
recognizer on their own subtrees. Meanwhile, \prop{certAffinePair} itself maps
recursively over all the children before combining their classifications.
The same nested subtree can therefore be classified repeatedly, once for each
nonlinear ancestor.\cite{certificates}

A claim that one call of \prop{certAffinePair} is linear in its input tree does
not imply that the surrounding evaluator is linear. The composition of an
$O(m)$ speculative recognizer with a recursively descending fallback needs its
own cost analysis.

\subsection{An exact family and visit count}

Define
\[
 e_0(x)=x,\qquad e_n(x)=x\bigl(1+e_{n-1}(x)\bigr),\quad n\ge1.
\]
This is the ordinary nested Horner form of $x+x^2+\cdots+x^{n+1}$; no expansion
into monomials is needed to construct it. Count one node for each operator,
constant occurrence, and variable occurrence in the expression tree. Then
\[
 N_n=4n+1.
\]
For $n\ge1$, $e_n$ is nonlinear. The intermediate sum $1+e_0$ is affine, while
$1+e_{n-1}$ is nonlinear for $n\ge2$.

Let $A_n$ be the number of \prop{certAffinePair} visits made by one recognition
attempt on $e_n$. Since the current recognizer maps all child expressions,
\[
 A_0=1,\qquad A_n=A_{n-1}+4=4n+1.
\]
Similarly, recognition of $1+e_{n-1}$ costs
\[
 B_n=A_{n-1}+2=4n-1.
\]

\begin{proposition}
For the inspected control flow, evaluating an enclosure of $e_n$ performs exactly
\[
 C_n=4n(n+1)
\]
affine-recognizer visits, provided that the source expression retains the stated
nested grammar and the evaluation reaches these helpers.
\end{proposition}
\begin{proof}
At $e_n$, the evaluator first makes the unsuccessful recognition attempt costing
$A_n$. It then evaluates the outer variable directly and visits $1+e_{n-1}$.
Recognition of that sum costs $B_n$. For $n\ge2$ it also fails, so the evaluator
continues into $e_{n-1}$. For $n=1$, the sum is affine and the recursion stops.
Thus $C_1=A_1+B_1=8$, and
\[
 C_n=C_{n-1}+A_n+B_n=C_{n-1}+8n.
\]
Summation gives $C_n=8\sum_{j=1}^n j=4n(n+1)$. Setting $C_0=0$ makes the same
recurrence valid at the first step as well.
\end{proof}

In terms of tree size,
\[
 C_n=\frac{N_n^2+2N_n-3}{4}.
\]
This is a quadratic number of structural recognition visits. It is not a
quadratic bit-complexity theorem for all rational arithmetic, nor a wall-clock
benchmark of \prop{InverseCertificate}.

\begin{table}[htbp]
\centering\small
\begin{tabular}{rrrr}
\toprule
Horner depth $n$ & Tree nodes $4n+1$ & Old recognizer visits & Prototype classifications \\
\midrule
1 & 5 & 8 & 4 \\
2 & 9 & 24 & 7 \\
4 & 17 & 80 & 13 \\
8 & 33 & 288 & 25 \\
16 & 65 & 1,088 & 49 \\
32 & 129 & 4,224 & 97 \\
64 & 257 & 16,640 & 193 \\
\bottomrule
\end{tabular}
\caption{Executed independent-model counts. The prototype shares the variable
leaf by node identity, hence $3n+1$ unique classifications; without that sharing,
one classification per tree node still gives a linear structural bound.}
\label{tab:counts}
\end{table}

\subsection{A semantics-preserving design}

A request-local bottom-up classification pass can compute each node's affine
pair, or a failed classification, once. The enclosure pass then consults those
results. An affine node still uses its exact range and one outward rounding. A
nonaffine node still uses the same recursive interval operations in the same
order. This is a targeted elimination of repeated recognition, not a global
symbolic-expression cache.

\begin{lstlisting}[language=Python]
def classify(node):
    if node.identity in local_classifications:
        return local_classifications[node.identity]
    child_pairs = [classify(child) for child in node.children]
    pair = combine_affine_pairs(node, child_pairs)
    local_classifications[node.identity] = pair  # pair or failure
    return pair

def enclose(node, interval, rounding_context):
    pair = local_classifications[node.identity]
    if pair is not failure:
        return round_outward(exact_affine_range(pair, interval))
    return original_interval_fallback(node, interval, rounding_context)
\end{lstlisting}

The implementation must not recursively hash or serialize the entire subtree on
every cache lookup, since doing so can reintroduce the same structural cost.
The supplied Python model uses live node identity and retains the map only for
one request. A Wolfram implementation could use a bounded annotated traversal or
stable local node identifiers; it needs its own measurement of actual allocation
and lookup behavior rather than assuming that association keys are free.

The interval-equivalence argument is a structural induction. For a rational
constant or the source symbol, both evaluators return the same enclosure. If a
composite node is affine, both classifiers compute the same exact pair and both
evaluators round the same exact range. If it is not affine, the induction
hypothesis gives equal child enclosures, and the same ordered arithmetic folds
give equal parent enclosures. No stronger interval theorem or new rounding rule
is introduced by retaining classifications.

The supplied model includes dyadic outward rounding, rather than comparing only
unrounded symbolic values. It checks identical intervals on seven Horner depths,
120 seeded random trees, and the huge exact-affine translation control at three
arithmetic precisions. Separate tests verify that the rounding helper encloses
2,982 rational values across the selected bit counts. These tests support the
prototype's stated grammar; they do not validate the entire Wolfram certificate
engine.

\subsection{Important limits of the result}

The exact count is conditional on the nested tree reaching \prop{certEnclose}
without prior flattening or polynomial expansion into another shape. Some public
certificate preparation paths simplify or expand expressions, and different
paths may evaluate a derivative or residual rather than the user's original tree.
A successful end-to-end native performance witness was not obtained here. The
finding is therefore a source-proved helper complexity defect with a reproducible
family, not an asserted timing regression for every public certificate call.

Likewise, one classification per node is an $O(N)$ structural-operation bound,
not necessarily $O(N)$ bytes or bit operations. Exact rational pairs can grow in
size, and storing pairs at several levels can retain more coefficient data than a
single traversal would. A production implementation should measure that cost and
bound any optional retained state. None of these qualifications changes the
quadratic repeated-visit count of the current control flow.

A focused native acceptance plan should instrument recognizer calls, compare
returned rational intervals exactly, retain affine translations and ordinary
nonlinear controls, and then benchmark a public path shown to preserve the
adversarial tree. Timing only a separately reimplemented classifier would not
establish a speedup of \prop{InverseCertificate}.

\section{N03: acceptance receipts can contradict their own source claims}
\label{sec:n03}

\subsection{What the checker already verifies}

The new checker is considerably more than a search for a passing count. Its
\prop{validate\_shard} function requires a Mathics runtime, completion,
source-stability and fresh-kernel flags, consistent selected/executed/succeeded
counts, zero failed/not-run counts, the expected source set and fingerprints,
a matching suite snapshot, and matching configured and effective iteration
budgets. Optional final fingerprints must agree with the initial map. Each case
must have the expected identity and group, a successful outcome and zero exit
code; raw protocol fields must agree with stored fields and the expected Mathics
version. The parsed printed actual and expected values must also agree. The
outer checker combines shard coverage against a Git-derived reference.\cite{checker}

Those are useful controls and should remain. The new finding is that some
individually plausible fields are never related to one another strongly enough.
The independently executed excerpt in this bundle contains the upstream
\prop{require}, \prop{canonical\_sources}, \prop{protocol\_fields}, and
\prop{validate\_shard} functions. It intentionally excludes the Git-discovery and
CLI layers. An additional supplied script can compare the four function ASTs
against a complete local checkout; that optional source-matching run was not
performed here. The tests therefore establish what the copied shard validator accepts,
not that an artificial bundle was run through every outer workflow command.

\subsection{Selected source versus fingerprinted source}

The checker first requires the selected basename to be the actual package entry,
then infers the layout from the immediate parent directory:
\begin{lstlisting}[language=Python]
parts = source.replace("\\", "/").split("/")
require(parts[-1] == ENTRY, "Unexpected package entry filename")
layout = "modular" if len(parts) >= 2 and parts[-2] == "Kernel" else "standalone"
\end{lstlisting}

It then independently canonicalizes keys in \prop{TestedSourcesSHA256} and checks
the resulting relative-name/digest dictionary against the expected dictionary.
For kernel files, canonicalization retains the final filename and the fact that
its parent is named \prop{Kernel}; it discards the earlier path components.
There is no subsequent requirement that the selected source path identify the
fingerprinted entry path.\cite{checker}

Start with an otherwise valid synthetic modular shard. Change only
\begin{lstlisting}[language=Python]
report["Source"] = "/NOT-FINGERPRINTED/src/Kernel/AsymptoticAnalysis.wl"
\end{lstlisting}
while leaving the fingerprint map under \prop{/checkout/}. The copied upstream
validator accepts the shard. The same mutation is accepted for a standalone
source when the basename remains \prop{AsymptoticAnalysis.wl}. In neither case
does the saved map contain the selected path.

The existing filename check is important: selecting
\path{/checkout/src/Kernel/InverseCertificates.wl} is \emph{already rejected}.
The bundled test retains that case as a negative control, not as a finding.
N03 concerns two differently located files with the \emph{same admitted entry
basename}, not arbitrary helper-module admission. Matching an entry name and
matching an independently supplied digest map do not establish that the two
records describe the same selected path.

\subsection{Recorded first drift versus stability flags}

The current review baseline records a runner repair that latches the first
observed mismatch even when later source content is restored. The point of that
record is monotonic invalidation: a later equal endpoint must not erase the fact
that a mismatch was already seen.\cite{w4}

The new checker requires a true source-stability flag and permits final fingerprints
to be absent or equal to the initial map. It does not inspect
\prop{FirstObservedSourceDriftSHA256}. Accordingly, an artificial receipt with
true stability flags, equal initial and final maps, and a non-null first-drift map
is accepted by the copied upstream validator. An empty dictionary in that
non-null field is also accepted; the candidate guard rejects both because the
current runner's no-drift sentinel is \prop{None}.

This is not another demonstration that endpoint fingerprints fail to prove
continuous immutability. It is narrower: when a record already contains explicit
evidence of observed drift, the validator should not accept mutually inconsistent
stability claims. Nor is it an allegation that the repository's real runner ever
emitted the particular contradictory fixtures. The tests deliberately construct
them to exercise the checker's consistency contract.

\begin{table}[htbp]
\small
\begin{tabularx}{\linewidth}{@{}Yp{0.19\linewidth}p{0.19\linewidth}@{}}
\toprule
Synthetic record & Upstream excerpt & Candidate guard \\
\midrule
Consistent modular positive control & Accepts & Accepts \\
Consistent standalone positive control & Accepts & Accepts \\
Selected entry absent from the fingerprint map & Accepts & Rejects \\
Fingerprint-listed helper selected as package entry (existing control) & Rejects & Rejects \\
First observed drift recorded despite equal endpoint maps & Accepts & Rejects \\
Existing failure controls: bad hash, incomplete flags, duplicate/empty cases,
missing protocol marker & Rejects & Rejects \\
\bottomrule
\end{tabularx}
\caption{Executed synthetic shard-validation behavior. No row is a real Mathics
execution receipt. The synthetic version strings explicitly identify the host as
artificial, and the fingerprints are dummy fixture data.}
\end{table}

\subsection{Candidate guard and its limits}

The supplied preflight guard adds two cross-field checks before delegating to
the original shard validator. The selected source's lexically normalized recorded
path must identify exactly one key in the fingerprint map. A present
\prop{FirstObservedSourceDriftSHA256} must be null. The existing entry-basename
check is preserved; it is not claimed as a new repair.

Path normalization is deliberately lexical. Windows drive and UNC syntax use
Windows case normalization; POSIX paths remain case-sensitive. The guard does
not resolve a receipt's paths on the reviewer's filesystem. Such resolution could
refer to unrelated files and would not establish the identity of a file on the
recorded execution host anyway. The current bundle does not claim symlink,
hard-link, filesystem-lock, or executed-byte attestation from this comparison.

A production integration should first check that the repository's real receipts
use a consistent selected-path spelling, then add fixtures for both layouts,
relative and absolute path spellings that are intentionally supported, and
legacy receipts with an absent drift field. If a supported entry format is
extended later, its admission rule and dependency contract should be updated
together. A supported new entry must remain linked to its own fingerprinted path.

These guards cannot make a maliciously fabricated receipt authentic. A party
able to rewrite every field can fabricate consistent fields as well. The value
is detection of accidental or adversarial \emph{internal contradictions} under
the checker's existing trust model. No signature, trusted-execution mechanism,
or cryptographic claim is proposed here.

\subsection{Why the finding is new rather than a runner replay}

The earlier wave-4 findings include live-source capture, source closure,
first-drift retention, mandatory package loading, and coverage aggregation. Their
recommended runner repairs should not be relabeled as new work. The present
counterexamples exercise a later checker that reads already saved records and
can accept contradictions those records make visible. They can be reproduced
without launching a child process, modifying a source file during execution,
loading Mathics, or bypassing a failed test. The additional obligation is at the
consumer of evidence, not its producer.\cite{w4,r36,checkercommit}

\section{Implementation priorities and focused development}
\label{sec:plan}

\subsection{First: close the concrete no-op provenance gap}

Apply the source-matched N01 change in the modular file, rebuild the standalone,
and run the six supplied native specifications together with the repository's
existing truncation-bound tests. Preserve the original quantitative bound and
conditions exactly on a no-op. Retain original-source objects without adding a
spurious discard history. The acceptance claim should name those selected tests
and the rebuilt artifact rather than imply complete package acceptance.

A useful immediate development check is contract dependency closure for this one
contract type: whenever a result stores
\prop{Type -> "TransportedThroughTruncation"}, access to the companion expression
used by its current statement should succeed. This is a narrow executable
invariant for a new metadata pair, not a proposal to redesign the entire result
schema or serialization system.

\subsection{Second: make the new evidence consumer fail on contradictions}

Integrate the N03 checks beside the existing shard checks, using the real receipt
layout as a compatibility control. Keep the current checks for fingerprints,
protocol consistency, execution budgets, and case coverage. A negative fixture
should state which single invariant it violates. This isolates regressions and
prevents a fixture rejected for an unrelated malformed field from falsely
appearing to test the desired guard.

The minimal test matrix has two valid entry layouts and two contradiction
families: an unfingerprinted selected path and a recorded first mismatch
contradicting the stability flag. A fingerprinted nonentry module remains an
existing rejection control. The supplied tests also
retain the old bad-hash, duplicate-case, empty-case, and protocol controls. The
outer CLI's complete-coverage check remains separate; the new guards should not
change how expected case coverage is computed.

\subsection{Third: remove repeated recognition without weakening certificates}

Implement a request-local classification pass or equivalent bounded annotated
traversal for N02. Before timing, count recognizer visits and compare all returned
rational endpoints exactly on the supported grammar. Preserve the huge affine
translation witness and valid nonlinear controls. Only then measure a public
certificate path whose expression shape is known to reach the changed helper.

The later development opportunity is to reuse a node's exact affine
classification across multiple interval evaluations within one certificate
request, because the pair depends on the rational expression and source symbol,
not on the interval. This extension needs an explicit request lifetime and must
not reuse a classification across a changed source expression, variable, or
unsupported grammar. It is not required for the initial repair, and no native
speedup for it is claimed. The prototype already shows the smaller first step:
one classification per node in a single enclosure call.

\subsection{What this plan intentionally does not repeat}

The existing register remains the place for the native-superset objective,
Mathics semantic compatibility, general refinement demand propagation,
multirate transseries, broad provenance compaction, and general proof-record
architecture. Those are substantial goals but not new findings here. The three
focused changes above can be reviewed and accepted independently without
reopening every older design decision.\cite{status,w3,w4}

\section{Validation, artifacts, and reproducibility}
\label{sec:validation}

\subsection{Executed independent checks}

The saved Python 3.13.5 run reports \textbf{34 passed} pytest cases. Fourteen cases
cover the exact interval model: seven depths, three affine-translation controls,
three parameterized rounding checks, and one seeded random-tree test. Fifteen
cases exercise the copied validator and candidate guards. Five cases exercise
the source-matched patcher's text handling and refusals. Some tests contain
multiple mathematical checks, but these are not inflated into an additional
package-test count.

Run from the extracted bundle root:
\begin{lstlisting}[language=bash]
python -m pytest -q tests
python code/generate_evidence.py
\end{lstlisting}

The second command writes \path{evidence/independent-results.json}. The saved text
version contains the same independent observations in a convenient readable
form. The reported affine quantities are operation counts, not native timing
measurements. The test duration printed by pytest is the duration of the Python
checks, not a performance result for the Wolfram package.

\subsection{Native reproduction and patch workflow}

Use a local checkout of the pinned revision and a fresh official kernel. The
supplied script accepts either the modular or the generated standalone entry:
\begin{lstlisting}[language=bash]
wolframscript -file code/probe_native.wls /absolute/path/AsymptoticAnalysis.wl
\end{lstlisting}

It records the selected path and expected review revision, but explicitly does
not attest the source bytes or Git state. Besides the N01 witness, it contains
optional helper instrumentation for N02. That instrumentation was not
successfully executed during this review. Its counters should be inspected
alongside the actual returned intervals before being used as native evidence.

To inspect and apply the N01 candidate locally:
\begin{lstlisting}[language=bash]
python code/apply_truncation_fix.py /path/to/checkout
python code/apply_truncation_fix.py /path/to/checkout --write
# From the checkout root, rebuild its standalone artifact:
python validation/build_standalone.py
\end{lstlisting}

After loading the chosen rebuilt entry, run
\prop{TestReport} on \path{code/NativeRegressions.wlt}. The baseline should fail
the companion-preservation assertions. A passing patched run is still an
acceptance task for the maintainer; this bundle does not contain such a native
receipt.

\subsection{Artifact map}

\begin{longtable}{@{}p{0.46\linewidth}p{0.47\linewidth}@{}}
\toprule
Path & Purpose and evidence status \\
\midrule
\endfirsthead
\toprule
Path & Purpose and evidence status \\
\midrule
\endhead
\path{article/review.tex}, \path{article/review.pdf} & Self-contained LaTeX source
and compiled report. \\
\path{code/apply_truncation_fix.py} & Opt-in, source-matched candidate N01 patch;
patcher fixtures executed, native package acceptance unrun. \\
\path{code/NativeRegressions.wlt} & Six post-fix native regression specifications;
unrun as a suite. \\
\path{code/probe_native.wls} & Reproduction script and native helper
instrumentation; script not accepted as an executed test run here. \\
\path{code/affine_model.py} & Independent bounded-grammar exact-rational and
dyadic interval model, including classification reuse. \\
\path{code/receipt_excerpt.py} & Copied upstream validation functions with source
attribution; Git and CLI layers intentionally omitted. \\
\path{code/receipt_guard.py}, \path{code/synthetic_receipts.py} & Candidate
consistency guard and explicitly artificial input records. \\
\path{code/check_upstream_excerpt.py} & Optional AST comparison against a supplied
local checkout; not run against a complete checkout here. \\
\path{code/generate_evidence.py} & Recomputes operation counts and synthetic
validation observations. \\
\path{tests/} & Executed Python tests and import configuration. \\
\path{evidence/native-witness.json} & Manual connector transcription, with
variable normalization, diagnostic and version-capture limitations. \\
\path{evidence/independent-results.json} & Executed independent-model and
synthetic-validator observations; not a package receipt. \\
\path{evidence/pytest-results.txt} & Saved 34-case Python test output. \\
\path{evidence/novelty-ledger.json} & Nearest prior findings, new mechanism, and
explicit nonclaims for each retained finding. \\
\bottomrule
\end{longtable}

\subsection{Limits that remain material}

No full repository suite was executed. No fresh Mathics package run was obtained.
The official-kernel evidence is a successful complete-load smoke and one public
truncation witness, not a combined MUnit acceptance run. The new interval
prototype covers a rational algebraic grammar, not transcendental enclosure
functions or all certificate domains. The receipt tests execute the specified
upstream shard-validation functions, not the full Git/CI acquisition workflow.
The novelty check uses the maintained consolidated inventory and targeted
original reports; it is not an independent line-by-line re-review of every old
article. These boundaries are reflected in the machine-readable evidence as well
as the prose.

\section{Conclusion}

The current snapshot contains three additional, sharply bounded opportunities
for improvement. Preserve the companion of a transported bound when truncation
is a no-op. Avoid repeatedly proving that the same nonlinear subtrees are not
affine. Reject source-entry and source-stability contradictions when consolidating
acceptance evidence. Each change preserves an existing useful mechanism rather
than replacing it: quantitative bound transport, exact affine cancellation, and
strict receipt checking.

The broader lesson is compositional. A correct bound formula is not a complete
metadata-copy policy. A linear recognizer is not necessarily a linear enclosing
algorithm. A set of individually checked receipt fields is not necessarily a
consistent receipt. The supplied witnesses, structural proof, independent
models, and targeted specifications make those three distinctions actionable
without recasting the repository's older backlog as new work.

\appendix
\section{Source locations and review anchors}

The function names below refer to the pinned repository files and provide
searchable locations without relying on line numbers of escaped connector responses. A moving branch should not be used to reconcile a mismatch.

\begin{longtable}{@{}p{0.12\linewidth}p{0.53\linewidth}p{0.26\linewidth}@{}}
\toprule
Finding & Pinned source & Relevant location \\
\midrule
\endfirsthead
\toprule
Finding & Pinned source & Relevant location \\
\midrule
\endhead
N01 & \path{src/Kernel/SeriesOperations.wl} & Bound-transport helper and following public truncation wrapper. \\
N01 & \path{src/Kernel/DirichletSpecialFunctions.wl} & Zeta construction, bound,
source coordinate, and metadata. \\
N02 & \path{src/Kernel/InverseCertificates.wl} & \prop{certAffinePair}, \prop{certAffineRange}, and \prop{certEnclose}. \\
N03 & \path{validation/check_mathics_acceptance.py} & \prop{canonical\_sources}, \prop{protocol\_fields}, and \prop{validate\_shard}. \\
Novelty & \path{docs/development/CODE_REVIEW_STATUS.md} & C20, C22, P08, and
maintained evidence boundaries. \\
Novelty & \path{external-reports/code-review/wave-2/code-review-14/article.tex} &
Original N02 affine-rounding and N05 quantitative-bound findings. \\
Novelty & \path{external-reports/code-review/wave-4/code-review-36/findings.json} &
Original runner/aggregation entries, especially N06. \\
\bottomrule
\end{longtable}

\section{Acceptance assertions in compact form}

For N01, the central postcondition is
\[
 \text{retained transported contract after a no-op}
 \ \Longrightarrow\ 
 \text{retained existing discard companion}.
\]
The finite expression, remainder, conditions, and quantitative upper bound should
be unchanged under the same no-op. This does not require retaining a signed
first-omitted-term certificate after a nontrivial truncation.

For N02, the first acceptance assertion is exact interval equality between old
and changed evaluators on the supported grammar. The second is a structural
counter bounded by a constant multiple of the input tree nodes, without recursive
whole-subtree key construction. A native wall-clock claim requires an additional
measured public workload and is not implied by either assertion.

For N03, existing entry-name admission, new path membership, and no-observed-drift are
conjoined with all existing shard checks. They are not alternatives to source
map equality, protocol checks, or complete case coverage. They establish
consistency under the existing recorder trust model, not authenticity of an
arbitrary submitted record.

\begin{thebibliography}{20}
\small
\bibitem{repo} Vladimir Reshetnikov and contributors.
\emph{Asymptotic}, reviewed revision \texttt{\pin}.
\url{https://github.com/VladimirReshetnikov/Asymptotic/tree/651f2029d0b2cd4da9e4dfdf1f4275a124d23b99}.
\bibitem{core} Repository source, \emph{AsymptoticAnalysis.wl}, canonical kernel.
\url{https://github.com/VladimirReshetnikov/Asymptotic/blob/651f2029d0b2cd4da9e4dfdf1f4275a124d23b99/src/Kernel/AsymptoticAnalysis.wl}.
\bibitem{operations} Repository source, \emph{SeriesOperations.wl}.
\url{https://github.com/VladimirReshetnikov/Asymptotic/blob/651f2029d0b2cd4da9e4dfdf1f4275a124d23b99/src/Kernel/SeriesOperations.wl}.
\bibitem{dirichlet} Repository source, \emph{DirichletSpecialFunctions.wl}.
\url{https://github.com/VladimirReshetnikov/Asymptotic/blob/651f2029d0b2cd4da9e4dfdf1f4275a124d23b99/src/Kernel/DirichletSpecialFunctions.wl}.
\bibitem{certificates} Repository source, \emph{InverseCertificates.wl}.
\url{https://github.com/VladimirReshetnikov/Asymptotic/blob/651f2029d0b2cd4da9e4dfdf1f4275a124d23b99/src/Kernel/InverseCertificates.wl}.
\bibitem{checker} Repository source, \emph{check\_mathics\_acceptance.py}.
\url{https://github.com/VladimirReshetnikov/Asymptotic/blob/651f2029d0b2cd4da9e4dfdf1f4275a124d23b99/validation/check_mathics_acceptance.py}.
\bibitem{checkercommit} Repository history, \emph{Verify and record complete
focused Mathics acceptance}, commit \texttt{e10ef1b71f5b55080c62f52c3ce9d3d6e039fc99}.
\url{https://github.com/VladimirReshetnikov/Asymptotic/commit/e10ef1b71f5b55080c62f52c3ce9d3d6e039fc99}.
\bibitem{reviews} Repository review corpus index, \emph{Code review reports}.
\url{https://github.com/VladimirReshetnikov/Asymptotic/blob/651f2029d0b2cd4da9e4dfdf1f4275a124d23b99/external-reports/code-review/README.md}.
\bibitem{status} Maintained \emph{Code review implementation status}.
\url{https://github.com/VladimirReshetnikov/Asymptotic/blob/651f2029d0b2cd4da9e4dfdf1f4275a124d23b99/docs/development/CODE_REVIEW_STATUS.md}.
\bibitem{w3} Maintained \emph{Wave 3 intake}.
\url{https://github.com/VladimirReshetnikov/Asymptotic/blob/651f2029d0b2cd4da9e4dfdf1f4275a124d23b99/docs/development/WAVE_3_INTAKE.md}.
\bibitem{w4} Maintained \emph{Wave 4 intake}.
\url{https://github.com/VladimirReshetnikov/Asymptotic/blob/651f2029d0b2cd4da9e4dfdf1f4275a124d23b99/docs/development/WAVE_4_INTAKE.md}.
\bibitem{r14} Original external report 14, \emph{Incremental Independent Audit of
AsymptoticInverse 1.8.0}, especially N02 and N05.
\url{https://github.com/VladimirReshetnikov/Asymptotic/blob/651f2029d0b2cd4da9e4dfdf1f4275a124d23b99/external-reports/code-review/wave-2/code-review-14/article.tex}.
\bibitem{r36} Original external report 36, finding ledger.
\url{https://github.com/VladimirReshetnikov/Asymptotic/blob/651f2029d0b2cd4da9e4dfdf1f4275a124d23b99/external-reports/code-review/wave-4/code-review-36/findings.json}.
\end{thebibliography}
\end{document}
